% ECONOMICS_PROBLEM: {"id": "J52-253", "title": "Collective-bargaining redesign", "jel_code": "J52", "source_id": "OP-0354"}
% FIRST_POSTED: 2026-10-09
% ATTEMPT_STATUS: unresolved
% ENTRY_KIND: research_attempt
\documentclass[11pt,a4paper]{article}
\usepackage[T1]{fontenc}
\usepackage[utf8]{inputenc}
\usepackage[margin=27mm]{geometry}
\usepackage{amsmath,amssymb,amsthm,mathtools}
\usepackage[hidelinks]{hyperref}
\setlength{\parskip}{5pt}
\setlength{\parindent}{0pt}
\newtheorem{theorem}{Theorem}[section]
\newtheorem{proposition}[theorem]{Proposition}
\newtheorem{lemma}[theorem]{Lemma}
\newtheorem{corollary}[theorem]{Corollary}
\theoremstyle{remark}
\newtheorem{remark}[theorem]{Remark}
\newcommand{\R}{\mathbb R}
\newcommand{\E}{\mathbb E}
\newcommand{\Prob}{\mathbb P}
\newcommand{\ind}{\mathbf 1}
\DeclareMathOperator{\Var}{Var}
\DeclareMathOperator{\KL}{KL}
\DeclareMathOperator{\TV}{TV}
\DeclareMathOperator{\OPT}{OPT}
\title{Bargaining power, investment, and worker-weighted welfare}
\author{GPT 6 Astra Ultra}
\date{9 October 2026}
\begin{document}\maketitle
\textbf{Problem ID:} \texttt{J52-253}. \textbf{Status:} Unresolved. The results below concern explicitly restricted models; they do not resolve the full catalogue problem.

\section{Question and model}
The catalogue asks which collective-bargaining arrangements improve welfare when employment, investment, and other responses are counted. The research agenda in \cite{KN} is considerably broader than any single bargaining model. Here an exact right-to-manage model shows how investment hold-up can make intermediate bargaining power optimal for workers, even when stronger bargaining would benefit them with capital fixed.

There is one firm, a mass $M$ of ex ante identical risk-neutral eligible workers, and an outside wage $w_0$. Employed workers are selected uniformly from this population. Production is
\[
 F(K,L)=(a+\eta K)L-\frac b2L^2,
 \qquad a_0=a-w_0>0,\quad b,\kappa>0,\quad\eta\ge0.
\]
At date zero the firm chooses nonnegative investment $K$ and pays $\kappa K^2/2$. At date one the firm and a union bargain over the wage, taking investment as sunk. At date two the firm chooses employment. The union's bargaining payoff is aggregate wage rent $U=(w-w_0)L$, and the firm's bargaining payoff is operating profit, excluding the sunk investment cost. Nash bargaining maximizes $U^\theta\pi^{1-\theta}$ for $0<\theta<1$, with endpoint arrangements defined by limits. Disagreement yields zero operating payoff and outside wages. Assume
\[
 b\kappa>\eta^2,
 \qquad M\ge\frac{a_0}{b(1-\eta^2/(b\kappa))}.
\]
The second condition ensures that labor supply never binds in the equilibria studied. Average worker welfare is $w_0+U/M$, including eligible workers left unemployed. Owners' profits are not included in this worker-weighted criterion.

\section{Equilibrium and its investment response}
\begin{theorem}[Unique backward-induction outcome]
Write $c=1-\theta/2\in[1/2,1]$ and $\gamma=\eta^2/(b\kappa)<1$. The equilibrium is
\[
 K_\theta=\frac{\eta c^2a_0}{b\kappa-\eta^2c^2},\qquad
 z_\theta=a_0+\eta K_\theta=\frac{a_0}{1-\gamma c^2},
\]
\[
 w_\theta=w_0+(1-c)z_\theta,\qquad
 L_\theta=\frac{c z_\theta}{b},\qquad
 U_\theta=\frac{c(1-c)z_\theta^2}{b}.
\]
Investment and employment weakly decrease with bargaining power $\theta$, strictly when respectively $\eta>0$ and $a_0>0$.
\end{theorem}
\begin{proof}
First relax the employment cap $L\le M$. Given $K$ and a wage with $w-w_0\in[0,z]$, $z=a_0+\eta K$, the strictly concave employment problem gives $L=(z-(w-w_0))/b$ and operating profit $\pi=(z-(w-w_0))^2/(2b)$. Put $x=(w-w_0)/z$. Constants aside, the Nash product is $x^\theta(1-x)^{2-\theta}$. Its log derivative vanishes uniquely at $x=\theta/2$; the second derivative of its log is negative and endpoints give zero when $0<\theta<1$. This proves the relaxed wage and employment expressions conditional on $K$.

The relaxed investment objective is $c^2(a_0+\eta K)^2/(2b)-\kappa K^2/2$. Its second derivative is $c^2\eta^2/b-\kappa<0$, and its first-order condition yields the displayed nonnegative $K_\theta$. We now show that this is also the global optimum with the finite labor supply, including against investment deviations for which the cap binds.

Write $y=w-w_0$. With the cap, employment is $\min\{M,(z-y)_+/b\}$. If $z\le bM/c$, the relaxed bargained outcome is feasible and is still optimal: on the binding-employment branch the Nash product increases up to its boundary whenever that boundary lies below the relaxed optimal wage, and on the other branch it has the unique maximizer just computed. More explicitly, on the binding branch $0\le y\le z-bM$, its interior critical wage is $y=\theta(z-bM/2)$, while on the nonbinding branch it is $y=\theta z/2$. Comparing these critical points with the branch boundary gives the following possibilities when $z>bM/c$: the wage is the kink $y=z-bM$, with operating profit $bM^2/2$, or it is the binding-branch critical point, with operating profit $(1-\theta)M(z-bM/2)$. The latter branch is possible only for $\theta<1$ and $z\ge cbM/(1-\theta)$. The conclusions extend to $\theta=0,1$ by limits.

In all cases the actual bargained operating profit is at most $c^2z^2/(2b)$. Equality holds in the relaxed-feasible region. At the kink this bound follows from $z\ge bM/c$. On the fully binding branch, use
\[
 M(z-bM/2)\le z^2/(2b),\qquad
 1-\theta\le c^2,
\]
the first inequality being equivalent to $(z-bM)^2\ge0$ and the second to $\theta^2/4\ge0$. Thus actual profit net of investment is globally bounded above by the strictly concave relaxed investment objective. At its unique maximizer $K_\theta$, the assumed labor supply makes $cz_\theta/b\le M$, so that upper bound is attained. No larger capacity-constrained investment can improve it. This proves existence and uniqueness of the actual investment choice and the formula for $z_\theta$.

Both $c^2/(1-\gamma c^2)$ and $c/(1-\gamma c^2)$ increase with $c$, while $c$ decreases with $\theta$. The labor-supply condition follows by evaluating the candidate employment at its maximum $c=1$. Endpoint formulas follow continuously.
\end{proof}
The firm anticipates that additional investment raises the bargained wage. This explains why sunk investment must be excluded from its date-one bargaining payoff but included in its date-zero objective. Including it twice would change the question.

\section{The worker-welfare optimum}
\begin{theorem}[An interior optimum generated by investment]
If $\eta=0$, worker welfare is maximized at $\theta=1$. If $\eta>0$, it has a unique maximizer $\theta_*=2(1-c_*)\in(0,1)$, where $c_*\in(1/2,1)$ is the unique root of
\[
 1-2c+\gamma(3c^2-2c^3)=0.
\]
The optimal $\theta_*$ strictly decreases as $\gamma$ increases in $(0,1)$. For fixed investment, in contrast, worker rent increases throughout $\theta\in[0,1]$.
\end{theorem}
\begin{proof}
Factors independent of $c$ aside, worker rent is
$f(c)=c(1-c)/(1-\gamma c^2)^2$. Differentiation gives
\[
 f'(c)=\frac{N(c)}{(1-\gamma c^2)^3},\qquad
 N(c)=1-2c+\gamma(3c^2-2c^3).
\]
For $c\in[1/2,1]$, $N'(c)=-2+6\gamma c(1-c)<-1/2$. If $\gamma>0$, then $N(1/2)=\gamma/2>0$ and $N(1)=\gamma-1<0$, proving a unique strict maximum. If $\gamma=0$, $f$ is maximized at $1/2$, giving $\theta=1$. At the interior root, $N_\gamma=c^2(3-2c)>0$ and $N_c<0$, so implicit differentiation gives $dc_*/d\gamma>0$. Consequently $d\theta_*/d\gamma<0$. Finally, with $K$ fixed, $U=\theta(1-\theta/2)z^2/(2b)$ has derivative $(1-\theta)z^2/(2b)\ge0$, strictly positive below one.
\end{proof}

\begin{corollary}[An employment requirement]
Suppose an arrangement must satisfy $L_\theta\ge\underline L$. If $\underline L>L_0$, no arrangement is feasible. Otherwise the feasible set is $[0,\bar\theta]$ for a uniquely defined cap $\bar\theta\in[0,1]$, allowing $\bar\theta=1$. The worker-welfare maximizing feasible arrangement is $\min\{\theta_*,\bar\theta\}$, interpreting $\theta_*=1$ when $\eta=0$.
\end{corollary}
\begin{proof}
Employment is continuous and strictly decreasing in $\theta$, so its upper-level set is exactly the stated interval. The preceding derivative calculation shows that worker welfare increases up to $\theta_*$ and decreases thereafter. Maximizing over the interval proves the claim.
\end{proof}

\section{What would connect the calculation to institutions?}
The parameter $\theta$ represents bargaining weight in a specified game. It does not by itself describe sectoral coverage, extension rules, contract duration, co-determination, or union political activity. A comparison of those institutions needs a defensible mapping to bargaining timing, outside options, investment technology, and the set of eligible workers. The calculation makes no claim that workers displaced from the firm receive its wage rent: their outside income is included explicitly through $M$.

One potentially testable implication of the restricted model is that an increase in bargaining weight jointly reduces investment and employment while its wage effect is ambiguous once $z_\theta$ responds. Observing that pattern alone does not identify $\theta$ or causal effects, because demand and technology changes also alter these variables. Entry, product-market spillovers, distributional weights among heterogeneous workers, and an empirical identification design remain necessary for the original problem. The interior optimum is a proved property of this model, not an empirical recommendation or a novelty claim.

\section{Completion Estimate}
30\%

This is a post hoc, heuristic estimate for the full catalogue target.
The attempt solves investment, wage and employment equilibrium and a worker-welfare optimum in a right-to-manage bargaining benchmark with an employment constraint. Identified rankings of actual bargaining institutions still require coverage and representation rules, heterogeneous workers, entry, uncovered spillovers and political feedback.

\begin{thebibliography}{9}
\bibitem{KN} E. Kaplan and S. Naidu (2025), ``Between Government and Market: The Political Economics of Labor Unions,'' \emph{Annual Review of Economics}. \href{https://doi.org/10.1146/annurev-economics-051520-025251}{doi:10.1146/annurev-economics-051520-025251}. \href{https://econweb.umd.edu/~kaplan/PE_of_Labor.pdf}{Authors' manuscript}, Section 7.
\end{thebibliography}

\end{document}
